\documentclass[11pt]{article}
\usepackage[margin=2.5cm]{geometry}
\usepackage{amsmath, amssymb, bm}
\usepackage{booktabs}
\usepackage{hyperref}

\newcommand{\body}{\mathrm{B}}
\newcommand{\ned}{\mathrm{NED}}
\newcommand{\flu}{\mathrm{FLU}}
\newcommand{\spf}{\mathrm{spf}}
\newcommand{\diag}{\operatorname{diag}}

\title{INDI Effectiveness Matrices: \\ NED (Indiflight) to FLU (Cybflight)}
\author{}
\date{}

\begin{document}
\maketitle

\section{Notation}

\begin{tabular}{ll}
    $N_u$ & Number of actuators (motors) \\
    $N_v$ & Number of pseudo-controls (6: $f_x, f_y, f_z, \tau_\phi, \tau_\theta, \tau_\psi$) \\
    $u_i \in [0,1]$ & Normalised motor command (force-proportional) \\
    $\bm{\omega}^\body$ & Body angular rate vector (rad/s) \\
    $\dot{\bm{\omega}}^\body$ & Body angular acceleration (rad/s$^2$) \\
    $\bm{a}_\spf^\body$ & Specific force in body frame (m/s$^2$), measured by accelerometer \\
    $\omega_{\mathrm{mot},i}$ & Motor $i$ angular speed (rad/s) \\
    $\dot{\omega}_{\mathrm{mot},i}$ & Motor $i$ angular acceleration (rad/s$^2$) \\
    $m$ & Vehicle mass (kg) \\
    $\bm{I}$ & Body inertia tensor (kg$\cdot$m$^2$), 3$\times$3 \\
    $T_i$ & Maximum thrust of motor $i$ (N) \\
    $\bm{r}_i = (p_{x,i}, p_{y,i}, 0)$ & Motor $i$ position in body frame (m) \\
    $s_i$ & Spin sign: $+1$ for CW-from-above, $-1$ for CCW \\
    $c_i$ & Torque coefficient (m): $|\tau_\mathrm{drag}| = c_i \cdot T_i$ \\
\end{tabular}

%% ============================================================================
\section{Indiflight Definitions (NED / FRD Frame)}
%% ============================================================================

Indiflight operates in NED/FRD: $x$ = forward, $y$ = right, $z$ = down.
Positive roll = right wing down, positive pitch = nose up, positive yaw = nose right.
Thrust acts in $-z$ (upward).

\subsection{G1: Static Effectiveness}

The $G_1$ matrix maps motor commands $\bm{u}$ to pseudo-controls
(specific forces and angular accelerations):

\begin{equation}
    \bm{\nu} = G_1 \, \bm{u}, \qquad
    G_1 \in \mathbb{R}^{6 \times N_u}
\end{equation}

Column $i$ of $G_1$ represents the effect of motor $i$ at full throttle ($u_i = 1$):

\begin{equation}\label{eq:g1_ned}
    G_1^{\ned}[:,i] = \begin{pmatrix}
        f_{x,i} \\[2pt] f_{y,i} \\[2pt] f_{z,i} \\[6pt]
        \alpha_{\phi,i} \\[2pt] \alpha_{\theta,i} \\[2pt] \alpha_{\psi,i}
    \end{pmatrix}
\end{equation}

For standard multirotors with all propellers producing thrust along body $-z_\ned$ (upward):

\begin{align}
    f_{x,i} &= 0 \label{eq:g1_fx_ned} \\
    f_{y,i} &= 0 \\
    f_{z,i} &= -\frac{T_i}{m} \qquad \text{(negative: thrust opposes $+z_\ned$ = down)} \label{eq:g1_fz_ned}
\end{align}

The angular acceleration rows are the torque divided by inertia.
In NED/FRD, the torque from thrust $\bm{F}_i = (0, 0, -T_i)$ at position
$\bm{r}_i = (p_{x,i}, p_{y,i}, 0)$ is:

\begin{equation}
    \bm{\tau}_i = \bm{r}_i \times \bm{F}_i + \bm{\tau}_{\mathrm{drag},i}
\end{equation}

Computing the cross product ($\bm{F}_i = (0, 0, -T_i)$ in NED):
\begin{align}
    \tau_{\phi,i}   &= p_{y,i} \cdot (-T_i) - 0 = -p_{y,i} \, T_i \\
    \tau_{\theta,i} &= 0 - p_{x,i} \cdot (-T_i) = +p_{x,i} \, T_i
\end{align}

For the yaw (drag reaction) torque: a CW-from-above propeller has positive
$\omega_z$ in NED ($z$ = down). The aerodynamic drag opposes the spin, so
drag acts on the propeller in $-z_\ned$.  By Newton's third law the body
reaction is also in $-z_\ned$ (the air pushes back on the body).
Therefore the yaw reaction torque is \emph{negative} for a CW motor ($s_i = +1$):
\begin{equation}
    \tau_{\psi,i} = -s_i \, c_i \, T_i
\end{equation}

The angular acceleration effectiveness:
\begin{equation}\label{eq:g1_alpha_ned}
    \begin{pmatrix} \alpha_{\phi,i} \\ \alpha_{\theta,i} \\ \alpha_{\psi,i} \end{pmatrix}
    = \bm{I}^{-1} \begin{pmatrix} \tau_{\phi,i} \\ \tau_{\theta,i} \\ \tau_{\psi,i} \end{pmatrix}
    = \bm{I}^{-1} \begin{pmatrix} -p_{y,i} \, T_i \\ +p_{x,i} \, T_i \\ -s_i \, c_i \, T_i \end{pmatrix}
\end{equation}

For diagonal inertia $\bm{I} = \diag(I_{xx}, I_{yy}, I_{zz})$:
\begin{align}
    \alpha_{\phi,i}   &= -p_{y,i} \, T_i \,/\, I_{xx} \\
    \alpha_{\theta,i} &= +p_{x,i} \, T_i \,/\, I_{yy} \\
    \alpha_{\psi,i}   &= -s_i \, c_i \, T_i \,/\, I_{zz}
\end{align}

\paragraph{Units.} Force rows: N/kg = m/s$^2$. Torque rows: Nm/(kg$\cdot$m$^2$) = rad/s$^2$.
Config scaling: \texttt{actG1\_fz} stored $\times 100$, \texttt{actG1\_roll/pitch/yaw} stored $\times 10$.

\subsection{G2: Rate-Dependent Effectiveness}

$G_2 \in \mathbb{R}^{3 \times N_u}$ captures the reaction torque from motor angular
acceleration. When motor $i$ accelerates ($\dot{\omega}_{\mathrm{mot},i} \neq 0$),
the propeller+rotor inertia $J_{\mathrm{prop},i}$ exerts a reaction torque on the body:

\begin{equation}
    \bm{\tau}_{\mathrm{react},i} = -J_{\mathrm{prop},i} \, \dot{\omega}_{\mathrm{mot},i} \, \hat{\bm{e}}_z
\end{equation}

For standard quads, only the yaw axis is affected:
\begin{equation}\label{eq:g2_ned}
    G_2^{\ned}[:,i] = \begin{pmatrix} 0 \\ 0 \\ G_{2,\psi,i} \end{pmatrix}
\end{equation}

The G2 values are configured directly (from system identification) and enter the
allocation via the G2 scaler:

\begin{equation}
    \text{G2\_scaler}_i = \frac{\omega_{\max,i}^2}{2 \, \tau_{\mathrm{mot},i}}
\end{equation}

where $\tau_{\mathrm{mot},i}$ is the motor time constant (s). The combined
effectiveness matrix used in the WLS allocation is:

\begin{equation}\label{eq:g1g2_combined}
    \widetilde{G}[j,i] = G_1[j,i] + \begin{cases}
        \dfrac{\text{G2\_scaler}_i}{\omega_{\mathrm{mot},i}} \, G_2[j-3,i]
            & \text{if } j > 2 \\[6pt]
        0 & \text{if } j \le 2
    \end{cases}
\end{equation}

\subsection{Pseudo-Control (INDI Law)}

The incremental pseudo-control error:
\begin{equation}\label{eq:dv_ned}
    \Delta\bm{\nu} = \begin{pmatrix}
        0 \\ 0 \\
        f_{z,\mathrm{sp}} - f_{z,\mathrm{meas}} \\
        \dot{\omega}_{\phi,\mathrm{sp}} - \dot{\omega}_{\phi,\mathrm{meas}} \\
        \dot{\omega}_{\theta,\mathrm{sp}} - \dot{\omega}_{\theta,\mathrm{meas}} \\
        \dot{\omega}_{\psi,\mathrm{sp}} - \dot{\omega}_{\psi,\mathrm{meas}}
    \end{pmatrix}
    + \begin{pmatrix} 0 \\ 0 \\ 0 \\
        \sum_i G_2[0,i] \, \dot{\omega}_{\mathrm{mot},i} \\
        \sum_i G_2[1,i] \, \dot{\omega}_{\mathrm{mot},i} \\
        \sum_i G_2[2,i] \, \dot{\omega}_{\mathrm{mot},i}
    \end{pmatrix}
\end{equation}

(Both terms multiplied by \texttt{doIndi} flag; when on ground, the controller
solves the non-incremental problem instead.)

%% ============================================================================
\section{Frame Transform: NED $\to$ FLU}
%% ============================================================================

Cybflight uses FLU: $x$ = forward, $y$ = left, $z$ = up.

\begin{equation}
    \begin{pmatrix} x \\ y \\ z \end{pmatrix}_{\!\flu}
    = \begin{pmatrix} 1 & 0 & 0 \\ 0 & -1 & 0 \\ 0 & 0 & -1 \end{pmatrix}
    \begin{pmatrix} x \\ y \\ z \end{pmatrix}_{\!\ned}
\end{equation}

\paragraph{Position.} $x$ unchanged, $y$ flips, $z$ flips:
\begin{equation}
    p_x^{\flu} = p_x^{\ned}, \qquad
    p_y^{\flu} = -p_y^{\ned}, \qquad
    p_z^{\flu} = -p_z^{\ned}
\end{equation}

\paragraph{Angular rates and torques.} Same transform (pseudovectors under proper rotation):
\begin{equation}
    \omega_\phi^{\flu} = \omega_\phi^{\ned}, \qquad
    \omega_\theta^{\flu} = -\omega_\theta^{\ned}, \qquad
    \omega_\psi^{\flu} = -\omega_\psi^{\ned}
\end{equation}

\paragraph{Rotor handedness label vs.\ yaw torque sign.}
The rotor handedness label $s_i$ is a geometric property of the propeller
(CW-from-above $= +1$, CCW $= -1$) and is the same in both frames.
However, the \emph{yaw torque} sign changes between frames because the
body $z$-axis flips:
\begin{alignat}{2}
    \tau_{\psi,i}^{\ned} &= -s_i \, c_i \, T_i &\qquad& \text{(CW $\to$ negative yaw in NED)} \\
    \tau_{\psi,i}^{\flu} &= +s_i \, c_i \, T_i && \text{(CW $\to$ positive yaw in FLU)}
\end{alignat}
This is consistent: $\tau_\psi^{\flu} = -\tau_\psi^{\ned}$ (z-axis flips).
Cybflight's \texttt{SpinDir::Cw = 1} directly gives the FLU sign.

%% ============================================================================
\section{Cybflight G1 Derivation (FLU Frame)}
%% ============================================================================

\subsection{From First Principles}

In FLU, thrust acts in $+z$ (upward). For motor $i$ with
$\bm{F}_i = (0, 0, +T_i)$ at $\bm{r}_i = (p_{x,i}^{\flu}, p_{y,i}^{\flu}, 0)$:

\begin{equation}
    \bm{\tau}_i^{\flu} = \bm{r}_i \times \bm{F}_i + \bm{\tau}_{\mathrm{drag},i}
\end{equation}

Cross product:
\begin{align}
    \tau_{\phi,i}^{\flu}   &= p_{y,i}^{\flu} \cdot T_i \label{eq:tau_roll_flu} \\
    \tau_{\theta,i}^{\flu} &= -p_{x,i}^{\flu} \cdot T_i \label{eq:tau_pitch_flu} \\
    \tau_{\psi,i}^{\flu}   &= s_i^{\flu} \cdot c_i \cdot T_i \label{eq:tau_yaw_flu}
\end{align}

These match \texttt{mixer.rs} lines 151--155 exactly.

The INDI G1 in FLU:
\begin{equation}\label{eq:g1_flu}
    \boxed{
    G_1^{\flu}[:,i] = \begin{pmatrix}
        0 \\[2pt]
        0 \\[2pt]
        T_i \,/\, m \\[6pt]
        \multicolumn{1}{c}{\bm{I}^{-1} \begin{pmatrix}
            p_{y,i}^{\flu} \, T_i \\
            -p_{x,i}^{\flu} \, T_i \\
            s_i^{\flu} \, c_i \, T_i
        \end{pmatrix}}
    \end{pmatrix}
    }
\end{equation}

For diagonal inertia:
\begin{align}
    G_1^{\flu}[0,i] &= 0 \\
    G_1^{\flu}[1,i] &= 0 \\
    G_1^{\flu}[2,i] &= +T_i / m & \text{(positive = upward in FLU)} \label{eq:g1_fz_flu} \\
    G_1^{\flu}[3,i] &= +p_{y,i}^{\flu} \, T_i / I_{xx} & \text{(roll)} \\
    G_1^{\flu}[4,i] &= -p_{x,i}^{\flu} \, T_i / I_{yy} & \text{(pitch)} \\
    G_1^{\flu}[5,i] &= +s_i^{\flu} \, c_i \, T_i / I_{zz} & \text{(yaw)}
\end{align}

\subsection{From Existing Cybflight G1 (Physical Units)}

Cybflight's \texttt{mixer.rs} already computes a 4-row G1 in physical units (FLU):
\begin{equation}
    G_{1,\mathrm{phys}}^{\flu}[:,i] = \begin{pmatrix}
        T_i \\
        p_{y,i}^{\flu} \, T_i \\
        -p_{x,i}^{\flu} \, T_i \\
        s_i^{\flu} \, c_i \, T_i
    \end{pmatrix}
    = \begin{pmatrix}
        \text{thrust (N)} \\
        \text{roll torque (N\,m)} \\
        \text{pitch torque (N\,m)} \\
        \text{yaw torque (N\,m)}
    \end{pmatrix}
\end{equation}

The INDI G1 is obtained by converting to acceleration space:
\begin{equation}\label{eq:g1_from_phys}
    \boxed{
    G_1^{\flu}[:,i] = \begin{pmatrix}
        0 \\[2pt]
        0 \\[2pt]
        G_{1,\mathrm{phys}}[0,i] \;/\; m \\[6pt]
        \multicolumn{1}{c}{\bm{I}^{-1} \begin{pmatrix}
            G_{1,\mathrm{phys}}[1,i] \\
            G_{1,\mathrm{phys}}[2,i] \\
            G_{1,\mathrm{phys}}[3,i]
        \end{pmatrix}}
    \end{pmatrix}
    }
\end{equation}

This handles non-diagonal inertia correctly via full $\bm{I}^{-1}$.

%% ============================================================================
\section{Cybflight G2 Derivation (FLU Frame)}
%% ============================================================================

G2 follows the same sign convention as the torque rows of G1.
For standard multirotors (vertical motors), only yaw is affected:

\begin{equation}\label{eq:g2_flu}
    \boxed{
    G_2^{\flu}[:,i] = \begin{pmatrix} 0 \\ 0 \\ G_{2,\psi,i}^{\flu} \end{pmatrix}
    }
\end{equation}

The yaw G2 sign equals the G1 yaw sign for the same motor:
\begin{itemize}
    \item CW motor ($s_i = +1$): speeding up $\Rightarrow$ reaction torque in $+z_{\flu}$
        $\Rightarrow$ positive yaw G2
    \item CCW motor ($s_i = -1$): speeding up $\Rightarrow$ reaction torque in $-z_{\flu}$
        $\Rightarrow$ negative yaw G2
\end{itemize}

G2 values come from system identification, not first principles.
The combined effectiveness matrix (same as Eq.~\ref{eq:g1g2_combined})
uses the G2 scaler to adapt based on operating point:

\begin{equation}
    \widetilde{G}^{\flu}[j,i] = G_1^{\flu}[j,i] +
    \begin{cases}
        \dfrac{\text{G2\_scaler}_i}{\omega_{\mathrm{mot},i}} \, G_2^{\flu}[j-3,i]
        & j > 2 \\[6pt]
        0 & j \le 2
    \end{cases}
\end{equation}

%% ============================================================================
\section{Numerical Verification}
%% ============================================================================

Using cybflight's SAKURAH743 quad parameters:
\begin{align*}
    m &= 0.55 \text{ kg}, \quad I_{xx} = 0.0025, \quad I_{yy} = 0.0021, \quad I_{zz} = 0.0043 \text{ kg\,m}^2 \\
    T_i &= 8.5 \text{ N}, \quad c_i = 0.022 \text{ m}
\end{align*}

\begin{center}
\begin{tabular}{lcccc}
    \toprule
    & M0 (RR, CW) & M1 (FR, CCW) & M2 (RL, CCW) & M3 (FL, CW) \\
    & $p=(-d,-d)$ & $p=(+d,-d)$ & $p=(-d,+d)$ & $p=(+d,+d)$ \\
    \midrule
    $f_z$ (m/s$^2$)           & $+15.45$ & $+15.45$ & $+15.45$ & $+15.45$ \\
    $\alpha_\phi$ (rad/s$^2$) & $-340.0$ & $-340.0$ & $+340.0$ & $+340.0$ \\
    $\alpha_\theta$ (rad/s$^2$)& $+303.6$ & $-303.6$ & $+303.6$ & $-303.6$ \\
    $\alpha_\psi$ (rad/s$^2$) & $+43.5$  & $-43.5$  & $-43.5$  & $+43.5$  \\
    \bottomrule
\end{tabular}
\end{center}

\noindent where $d_x = 0.075$ m, $d_y = 0.1$ m. Physical checks:
\begin{itemize}
    \item All $f_z > 0$: thrust is upward in FLU. \checkmark
    \item M0 (rear-right): $\alpha_\phi < 0$ (right motor $\Rightarrow$ left-side-up roll). \checkmark
    \item M0 (rear): $\alpha_\theta > 0$. In FLU, positive $\tau_y$ rotates $+x$ toward
        $-z$ via $\bm{\omega} \times \bm{r}$, i.e.\ nose down. Rear motor thrusting
        up pitches nose down. \checkmark
    \item M0 (CW): $\alpha_\psi > 0$ (CW reaction $= +z_{\flu}$). \checkmark
    \item Yaw signs: CW motors (M0, M3) positive, CCW (M1, M2) negative. \checkmark
    \item Roll: right motors (M0, M1, $p_y < 0$) negative, left (M2, M3) positive. \checkmark
    \item Pitch: rear motors (M0, M2, $p_x < 0$) positive, front (M1, M3) negative. \checkmark
\end{itemize}

\paragraph{Cross-check with NED.} Applying the frame transform to the FLU values:
\begin{align*}
    G_1^{\ned}[2,\text{M0}] &= -G_1^{\flu}[2,\text{M0}] = -15.45 & \text{($-z$ = upward in NED)} \\
    G_1^{\ned}[3,\text{M0}] &= +G_1^{\flu}[3,\text{M0}] = -340.0 & \text{(roll, $x$ unchanged)} \\
    G_1^{\ned}[4,\text{M0}] &= -G_1^{\flu}[4,\text{M0}] = -303.6 & \text{(pitch, $y$ flips)} \\
    G_1^{\ned}[5,\text{M0}] &= -G_1^{\flu}[5,\text{M0}] = -43.5  & \text{(yaw, $z$ flips)}
\end{align*}

In NED: CW motor yaw is negative (drag reaction in $-z_{\ned}$), rear motor pitch is
negative (nose-down in NED sign convention). All consistent with indiflight. \checkmark

\end{document}
